%====================================
% ANHANG: Dokumentation der KI-Prompts
%====================================
\section*{Anhang: Dokumentation der KI-Prompts}
\addcontentsline{toc}{section}{Anhang: KI-Prompts}

Im Folgenden sind die zentralen Prompts aufgeführt, die im Rahmen der Erstellung dieses Exposés gegenüber KI-Sprachmodellen (Claude) verwendet wurden. Die Antworten der KI sind nicht dokumentiert, da nur die vom Verfasser eigenständig geprüften und angepassten Inhalte in die Arbeit eingeflossen sind.

\begin{enumerate}
  \item \glqq{}Hast du Vorschläge für Suchbegriffe zu ERP-Erfolgsfaktoren und Change Management im Mittelstand?\grqq{}
  
  \item \glqq{}Welche booleschen Suchstrings könnten für eine systematische Literaturrecherche zu ERP-Erfolgsfaktoren und Change Management geeignet sein?\grqq{}
  
  \item \glqq{}Welche Datenbanken eignen sich für eine systematische Literaturrecherche in der Wirtschaftsinformatik?\grqq{}
  
  \item \glqq{}Welche Kapitelstruktur wäre für eine Bachelorarbeit zu ERP-Change-Management im Mittelstand sinnvoll?\grqq{}
  
  \item \glqq{}Kannst du mir eine Alternative für diesen Satz vorschlagen? Er wirkt noch zu umständlich.\grqq{}
  
  \item \glqq{}Ist die Formulierung der Forschungsfrage präzise genug oder sollte ich sie eingrenzen?\grqq{}
  
  \item \glqq{}Passt der FOM-Fußnotenstil hier so?\grqq{}
\end{enumerate}
\newpage
